\chapter{Just Enough Complex Analysis}

Chapter~1 treated $z$ as a formal symbol, manipulated according to algebraic rules
with no thought of convergence or magnitude.
The real power of analytic combinatorics comes from promoting $z$ to a genuine
complex variable and extracting the asymptotic growth of the coefficients
$[z^n] A(z)$ from the analytic structure of the function $A$.
This chapter develops the minimum complex analysis required to make that move
rigorous and computational.
The reader who has seen complex numbers but not complex functions will find
everything needed here; the reader who has taken a full complex analysis course
may skim and use this chapter as a reference for notation.
The standard treatment of all topics below is \cite{FlajoletSedgewick2009}, Part~A.

\section{Complex Numbers and Convergence}

Recall that $\mathbb{C} = \{x + iy : x, y \in \mathbb{R}\}$ with the convention
$i^2 = -1$.
The \emph{modulus} of $z = x + iy$ is $|z| = \sqrt{x^2 + y^2}$, and the
\emph{argument} $\arg z$ is the angle $\theta \in (-\pi, \pi]$ such that
$x = |z|\cos\theta$, $y = |z|\sin\theta$.
Euler's formula
\[
  e^{i\theta} = \cos\theta + i\sin\theta
\]
allows every nonzero complex number to be written in \emph{polar form}
$z = r e^{i\theta}$ with $r = |z| > 0$ and $\theta = \arg z$.
Multiplication becomes transparent in polar form: $r_1 e^{i\theta_1} \cdot
r_2 e^{i\theta_2} = r_1 r_2 \, e^{i(\theta_1+\theta_2)}$, so moduli multiply
and arguments add.

Convergence in $\mathbb{C}$ is measured by the modulus: a sequence $(z_k)$
converges to $L$ if $|z_k - L| \to 0$.
A series $\sum_k w_k$ \emph{converges absolutely} if $\sum_k |w_k| < \infty$,
and absolute convergence implies convergence.
The geometric series identity
\[
  \sum_{k=0}^{N} w^k = \frac{1 - w^{N+1}}{1 - w}, \quad w \neq 1,
\]
shows that $\sum_{k=0}^{N} w^k \to 1/(1-w)$ as $N \to \infty$ whenever
$|w| < 1$, since $|w|^{N+1} \to 0$.
We will use this repeatedly.

\section{Power Series and Radius of Convergence}

Let $A(z) = \sum_{n \ge 0} a_n z^n$ be a formal power series with coefficients
$a_n \in \mathbb{C}$.
Treating $z$ as a complex variable, the set of $z$ for which this series converges
absolutely is a disk.

\begin{theorem}[Cauchy--Hadamard]
  The power series $A(z) = \sum_{n \ge 0} a_n z^n$ converges absolutely for every
  $z$ with $|z| < R$ and diverges for every $z$ with $|z| > R$, where the
  \emph{radius of convergence} $R$ is given by
  \[
    \frac{1}{R} = \limsup_{n \to \infty} |a_n|^{1/n}.
  \]
\end{theorem}

The formula follows from comparing $|a_n z^n|^{1/n} = |a_n|^{1/n} |z|$
with the ratio test: absolute convergence holds when this quantity is eventually
less than $1$, i.e., when $|z| < 1/\limsup |a_n|^{1/n}$.

The combinatorial import of Cauchy--Hadamard is immediate: the radius of
convergence captures the \emph{exponential growth rate} of $(a_n)$.
If $R > 0$ is finite, then $a_n$ grows (roughly) like $R^{-n}$, in the sense that
$|a_n|^{1/n} \to R^{-1}$.
More precisely, $a_n = R^{-n} \cdot \theta(n)$ where $\theta(n)$ is
\emph{subexponential}: $|\theta(n)|^{1/n} \to 1$.
Determining the fine structure of $\theta(n)$ is precisely what asymptotic
analysis does.

\begin{example}
  For $B(z) = 1/(1-2z) = \sum_{n \ge 0} 2^n z^n$, we have $b_n = 2^n$ and
  $|b_n|^{1/n} = 2$ for all $n$, so $R = 1/2$.
  The singularity of $B$ at $z = 1/2$ is directly visible as a pole,
  and the exponential growth rate $2^n = (1/2)^{-n}$ agrees with $R^{-n}$.

  For the Fibonacci generating function $F(z) = z/(1 - z - z^2)$,
  the poles are at the roots of $1 - z - z^2 = 0$, namely $z = (-1 \pm \sqrt{5})/2$.
  The root of smallest modulus is $1/\varphi$ where $\varphi = (1+\sqrt{5})/2
  \approx 1.618$ is the golden ratio, so $R = 1/\varphi$ and
  $f_n \sim C \varphi^n$ for some explicit constant $C$.
\end{example}

\section{Analytic Functions}

A function $f: U \to \mathbb{C}$ defined on an open set $U \subseteq \mathbb{C}$
is \emph{holomorphic} on $U$ (also called \emph{analytic}) if it is
complex-differentiable at every point of $U$.
A fundamental and non-trivial theorem of complex analysis states that the
following conditions are equivalent for a function on an open set:
(1) $f$ is complex-differentiable (the limit $\lim_{h \to 0}(f(z+h)-f(z))/h$
exists in $\mathbb{C}$);
(2) $f$ is infinitely differentiable and satisfies the Cauchy--Riemann equations
$\partial_x u = \partial_y v$, $\partial_y u = -\partial_x v$ where $f = u + iv$;
(3) $f$ is locally representable by a convergent power series around every point.
We will take this equivalence as given.

The key example: if $A(z) = \sum_{n \ge 0} a_n z^n$ has radius of convergence
$R > 0$, then $A$ is holomorphic on the open disk $|z| < R$.
This follows because power series can be differentiated term by term within
their disk of convergence, and the resulting series again converges.
Conversely, a holomorphic function on $|z| < R$ is completely determined by
its power series at the origin.
The formal generating functions of Chapter~1, as soon as their coefficients
grow at most geometrically, are genuine holomorphic functions on a nontrivial disk.

\section{The Geometric Series as Prototype}

The function $f(z) = 1/(1-z)$ is the model example.
For $|z| < 1$ the series $\sum_{n \ge 0} z^n$ converges absolutely, and
multiplying by $(1-z)$ collapses the partial sums to $1 - z^{N+1} \to 1$, giving
\[
  \frac{1}{1-z} = \sum_{n \ge 0} z^n, \quad |z| < 1.
\]
The function $1/(1-z)$ is holomorphic on all of $\mathbb{C} \setminus \{1\}$;
it extends analytically well beyond the disk of convergence of the series,
but the series itself is unable to see past the boundary $|z|=1$.
The point $z = 1$ is a \emph{simple pole}: $f(z) = 1/(1-z)$ blows up like
$(1-z)^{-1}$ near $z = 1$, which is the simplest possible singularity.

The lesson is structural: the singularity at $z=1$ is what \emph{prevents}
the radius of convergence from exceeding $1$, because the series cannot
represent a function that blows up.
In general, the radius of convergence of a power series at a point $z_0$
equals the distance from $z_0$ to the nearest singularity of the function
that the series represents.
Every analytic-combinatorics computation ultimately reduces to locating
the nearest singularity of a generating function.

\section{Cauchy's Integral Formula for Coefficients}

The coefficients of a power series can be recovered by a contour integral.
Let $f$ be holomorphic on $|z| < R$ and let $0 < r < R$.
Then
\[
  a_n = [z^n] f(z) = \frac{1}{2\pi i} \oint_{|z| = r} \frac{f(z)}{z^{n+1}}\, dz,
\]
where the contour integral runs counterclockwise around the circle of radius $r$.
This follows from Cauchy's integral formula for derivatives: for a holomorphic
$f$ on an open disk,
\[
  f^{(n)}(0) = \frac{n!}{2\pi i} \oint_{|z|=r} \frac{f(z)}{z^{n+1}}\, dz,
\]
combined with Taylor's theorem $f(z) = \sum_n a_n z^n$ implying
$f^{(n)}(0) = n!\, a_n$.
Dividing both sides by $n!$ gives the coefficient formula.

The crucial observation for asymptotics is that the contour $|z| = r$ can be
\emph{deformed outward}, provided it remains in the region where $f$ is
holomorphic.
On a circle of radius $r > 1$ (say), the integrand $|f(z)/z^{n+1}|$ is bounded
above by $\max_{|z|=r} |f(z)| \cdot r^{-(n+1)}$, which decays exponentially in
$n$ at rate $r^{-1}$.
A larger $r$ gives a tighter upper bound, and hence a sharper estimate on $a_n$.
The contour can be pushed outward as far as possible without crossing a
singularity of $f$; the singularity forms an impassable barrier.
It is at the nearest singularity---the one that stops the contour---that
the dominant contribution to $a_n$ resides.
This is the geometric heart of singularity analysis: singularities are not
obstacles to computation but rather the \emph{source} of asymptotic information.

\section{Classifying Singularities}

The type of a singularity determines the subexponential correction to the
exponential growth rate.
Let $\rho > 0$ be the location of the dominant singularity (so $|\rho|$ equals
the radius of convergence).
The three types that appear most frequently in combinatorics are the following.

\textbf{Poles.}
If $f(z) = 1/(1 - z/\rho)^k$ for a positive integer $k$, then by the
generalized binomial theorem,
\[
  \frac{1}{(1 - z/\rho)^k} = \sum_{n \ge 0} \binom{n + k - 1}{k - 1} \rho^{-n} z^n,
\]
so $[z^n] f = \binom{n+k-1}{k-1} \rho^{-n} \sim \frac{n^{k-1}}{(k-1)!} \rho^{-n}$.
The subexponential factor is a polynomial of degree $k-1$ in $n$.
A simple pole ($k=1$) contributes a pure exponential; a double pole ($k=2$)
contributes a linear factor $n \cdot \rho^{-n}$; and so on.
Most rational generating functions have poles, and they arise naturally whenever
a linear recurrence has a unique dominant root.

\textbf{Square-root branch points.}
If $f(z) = \sqrt{1 - z/\rho} = (1 - z/\rho)^{1/2}$, then the generalized
binomial theorem gives $f(z) = \sum_{n \ge 0} \binom{1/2}{n} (-z/\rho)^n$.
Computing the coefficient $\binom{1/2}{n}(-1)^n$ by Stirling's formula yields
\[
  [z^n]\sqrt{1 - z/\rho} = -\frac{1}{2\sqrt{\pi}}\, n^{-3/2}\, \rho^{-n}\bigl(1 + O(1/n)\bigr).
\]
The exponent $-3/2$ in the polynomial correction $n^{-3/2}$ is the
\emph{fingerprint} of a square-root singularity.
This type arises pervasively in combinatorics because the kernel method and
quadratic equations produce square roots; trees, lattice paths, and
many self-referential structures lead to generating functions with
$\sqrt{1 - z/\rho}$ as a factor.
We will return to this in Chapter~5.

\textbf{Logarithmic singularities.}
If $f(z) = \log(1/(1 - z/\rho))$, then $f(z) = \sum_{n \ge 1} \rho^{-n} z^n / n$,
so $[z^n] f = \rho^{-n}/n$.
The correction factor is $1/n$, slower to decay than any negative power with
exponent greater than $-1$ but faster than a constant.
Logarithms appear, for instance, in the generating functions for labelled
structures involving cycles, and in the analysis of certain algorithms.
They are less common than poles or branch points but arise naturally whenever
the symbolic construction includes $\SET$ or $\CYC$ operators.

\textbf{Essential singularities.}
Functions such as $e^{1/(1-z)}$ have an essential singularity at $z = 1$:
they blow up in a highly oscillatory, non-power-law fashion governed by
Picard's theorem.
The elementary transfer-theorem machinery does not apply at essential
singularities, and special methods are required.
Such singularities are rare in the combinatorial generating functions
considered in this book, and we will not develop the necessary tools here.

\section{The Big Picture}

The moral of this chapter is simple to state and powerful to apply.
If $A(z)$ is a generating function whose nearest singularity to the origin
lies at $|z| = \rho$, then $a_n$ grows exponentially like $\rho^{-n}$.
That much is determined purely by the \emph{location} of the singularity.
The \emph{type} of the singularity---pole of order $k$, square-root branch
point, logarithm---controls the subexponential factor $n^{\alpha}$ for some
$\alpha$ determined by the local singular behavior.
Flajolet and Odlyzko's transfer theorem, developed in Chapter~4, makes
this dictionary precise and gives explicit asymptotic expansions to any
desired order of precision \cite{FlajoletSedgewick2009}.

All the complex analysis we use in this book reduces to three steps:
(i) find a useful closed form for the generating function $A(z)$, using the
symbolic method from Chapter~1; (ii) locate the singularities on the boundary
of the disk of convergence; (iii) read the asymptotics off their type using
the classification above.
The deeper machinery of complex analysis---contour deformation through branch
cuts, analytic continuation to Riemann surfaces, the residue theorem in full
generality, the saddle-point method---belongs to a dedicated complex-analysis
course and will appear in this book only when no simpler approach suffices.
The examples of Chapter~3 will make the three-step program concrete before
we prove the transfer theorem in Chapter~4.
